\section{Conceptualizing baselines}
\label{app:baselines}

We summarize the baseline methods used in our experiments with an emphasis on the generation/sampling mechanism for each method. 

\paragraph{GraphEBM} ~\citep{liu2021graphebm} models molecular graphs using an energy-based model that assigns a scalar energy $E_\theta(x)$ to each graph $x$ via a graph neural network.
Unlike transport- or flow-based approaches, GraphEBM does not learn an explicit transformation from a source distribution to the data distribution. Instead, generation depends exclusively on Langevin dynamics sampling within a continuous embedding space.


\paragraph{DeFoG}~\citep{qinmadeira2024defog} adapts discrete flow matching to graphs by pairing an explicit noising path (via linear interpolation toward a simple base distribution) with a learned \gls{ctmc} denoiser: a network predicts the clean-graph posterior from an intermediate noisy graph, which in turn defines the time-dependent rate matrices used for generation.
Sampling starts from the base distribution and simulates the resulting \gls{ctmc} toward the data distribution, with the sampling step schedule largely selectable at inference time.

\paragraph{\gls{vfm}}~\citep{eijkelboom2024vfm} casts flow matching as variational inference over trajectory endpoints (the “posterior probability path”), yielding a KL-based objective that for categorical graph variables reduces to a per-component cross-entropy.
In contrast to DeFoG’s stochastic \gls{ctmc} jump dynamics, \gls{vfm} targets a deterministic continuous flow (vector field) on the probability simplex and generates by integrating this flow from the base distribution, then discretizing/sampling from the final categorical probabilities.


\paragraph{G-\gls{vfm}}~\citep{eijkelboom2025controlledvfm} extends \gls{vfm} with equivariant conditioning mechanisms to
support controllable generation.
